\section{Conclusion}
\label{sec:conclusion}

We have presented the largest cross-factorial benchmark to date for evaluating foundation model embeddings as featurisers for probabilistic surrogates in materials property prediction.
Across 8{,}640 cross-validated experiments spanning 4~featurisers, 3~surrogates, 3~PCA levels, 3~training sizes, and 3~materials property datasets---elastic tensor (8~properties), dielectric constant (4~properties), and phonon dielectric (3~properties)---six key findings emerge.

\textbf{1.\ ORB embeddings are the universal winner.}
ORB outperforms all alternatives on every dataset and every surrogate model.
On elastic properties, ORB\,+\,GP achieves $\bar{R}^2 = 0.803$ and $\bar{\rho} = 0.842$, compared to 0.731/0.820 for MACE, 0.734/0.811 for UMA, and 0.445/0.688 for SOAP.
On dielectric properties, ORB\,+\,DGP leads with $\bar{R}^2 = 0.336$ and $\bar{\rho} = 0.778$.
On phonon properties, ORB\,+\,GP achieves $\bar{R}^2 = 0.303$ and $\bar{\rho} = 0.824$.
This dominance stems from ORB's large and diverse pre-training corpus ($\sim$120M configurations) and its graph attention architecture optimised for periodic crystals.

\textbf{2.\ Ranking accuracy far exceeds regression accuracy.}
Across all three datasets, Spearman rank correlations substantially exceed \Rtwo{} values.
The gap widens with property difficulty: modest on elastic data ($\bar{\rho} = 0.842$ vs.\ $\bar{R}^2 = 0.803$), large on dielectric data ($\bar{\rho} = 0.778$ vs.\ $\bar{R}^2 = 0.336$), and dramatic on phonon data ($\bar{\rho} = 0.824$ vs.\ $\bar{R}^2 = 0.303$).
The extreme case of eps\_electronic ($\rho = 0.930$, $R^2 = -0.001$) demonstrates that surrogates can correctly order 93\% of material pairs even when absolute predictions fail entirely.
For Bayesian optimisation---where acquisition functions depend on relative ranking, not absolute prediction---this means our surrogates are useful across \emph{all three property domains}, not just elastic.
We recommend that the community adopt Spearman correlation as a primary metric for BO surrogates.

\textbf{3.\ The GP is the most reliable surrogate; the DGP is the most powerful when properly configured.}
The standard GP offers the best balance of accuracy, stability, and computational efficiency, serving as a robust default across all datasets and PCA settings.
The DGP achieves the single best \Rtwo{} on both the elastic dataset ($\bar{R}^2 = 0.807$ at PCA\,=\,25) and the dielectric dataset ($\bar{R}^2 = 0.336$ at PCA\,=\,10), demonstrating its ability to capture non-linear structure--property relationships that the linear GP cannot.
The key practical consideration is PCA dimensionality: the DGP performs strongly at PCA\,$\leq$\,25 but requires careful configuration at higher dimensions.

\textbf{4.\ Property predictability follows a physical gradient.}
Properties that depend on local bonding geometry (bulk modulus, $R^2 > 0.91$; phonon DOS peak, $R^2 = 0.911$) are substantially easier than properties that require electronic structure information (band gap, $R^2 = 0.852$; refractive index, $R^2 = 0.552$) or involve collective many-body responses (polycrystalline dielectric constants, $R^2 < 0$; phonon-derived dielectric constants, $R^2 \approx -0.001$).
This gradient maps directly onto the ``information distance'' between crystal geometry---what the embeddings encode---and the target physical phenomenon.

\textbf{5.\ Foundation model embeddings are remarkably data-efficient.}
ORB\,+\,GP at $\ntrain = 100$ surpasses SOAP\,+\,GP at $\ntrain = 500$ on the elastic dataset ($\bar{R}^2 = 0.535$ vs.\ $0.445$), providing a head start equivalent to $\sim$400 additional DFT calculations.
For BO campaigns where each evaluation costs hours of compute, this efficiency translates directly to shorter discovery timelines.

\textbf{6.\ Multi-task GPs help selectively.}
The MTGP shows no advantage over independent GPs on the elastic dataset ($T = 8$ tasks), but matches or exceeds the GP on the phonon dataset ($T = 3$; $\bar{\rho} = 0.830$ vs.\ 0.824) and improves over GP for weaker featurisers on the dielectric dataset (SOAP\,+\,MTGP: $\bar{R}^2 = 0.177$ vs.\ SOAP\,+\,GP: 0.132).
Multi-task learning is most beneficial when the number of tasks is small, properties share a common physical origin, and per-task prediction is weak.

\emph{Our updated practical recommendation:} Use ORB\,+\,GP\,+\,PCA\,=\,50 as the default stack for materials BO campaigns.
Upgrade to DGP at PCA\,$\in \{10, 25\}$ for non-linear properties with sufficient data ($\ntrain \geq 300$).
Evaluate surrogates by Spearman $\rho$, not only \Rtwo{}, especially for heavy-tailed or electronic properties.
For heavy-tailed targets, apply log-transforms before training to close the Spearman--\Rtwo{} gap.

\textbf{Limitations.}
Our benchmark uses three datasets with up to 500~training samples each, all comprising ordered crystalline structures.
We do not evaluate the full BO loop (acquisition function optimisation and sequential candidate selection), uncertainty calibration (reliability diagrams, ENCE), or the effect of fine-tuning foundation model embeddings on the target property.
The DGP's sensitivity to PCA dimension suggests that its performance ceiling may be higher with adaptive variational inference schemes.
At larger sample sizes ($\ntrain > 1{,}000$), the surrogate ranking---particularly for MTGP and DGP---may shift.

\textbf{Future work.}
Six directions emerge from the cross-dataset analysis:
(i)~integration into a full BO loop to measure how Spearman-based surrogate quality translates to closed-loop discovery efficiency;
(ii)~extension to disordered alloy, amorphous, and MOF systems where the featuriser ranking may differ;
(iii)~fine-tuning of foundation model embeddings on the target property, which could close the gap between ORB and MACE/UMA or between geometric embeddings and electronic-structure-aware models for dielectric properties;
(iv)~rigorous uncertainty calibration analysis for each surrogate--featuriser combination, which is critical for BO acquisition functions;
(v)~target-variable preprocessing (log-transforms, robust scaling) to close the Spearman--\Rtwo{} gap on heavy-tailed properties, potentially converting ``unpredictable'' properties like eps\_electronic ($\rho = 0.930$, $R^2 = -0.001$) into well-predicted ones;
and (vi)~active learning experiments where the surrogate actively selects the next training points, testing whether the DGP's richer posterior improves exploration--exploitation balance relative to the GP.

The complete pipeline, aggregated results, and analysis scripts are available as open-source software to serve as a community benchmark for future featuriser and surrogate development.
